% Abstract and index terms

\begin{abstract}
DataOps and MLOps are commonly adopted side by side, yet the tools supporting them remain fragmented and the two disciplines are usually run by separate teams. A team that wants both must either assemble the integration itself, which takes time, or subscribe to managed services, which carries recurring cost and ties the platform to one provider. This paper develops an integrated DataOps and MLOps platform architecture on Kubernetes built from open source tools. Design Science Research Methodology guides the work from problem identification through design, GitOps-based implementation, and evaluation. The architecture comprises nine platform layers derived from the convergence of five sources that used different methods rather than from majority adoption, and components for each layer are selected through scored comparison against four criteria. The design joins the two disciplines at three points, namely a shared feature contract, an open lineage protocol feeding one metadata catalog, and one declarative control plane reconciled from a single Git repository. A crypto asset market use case verifies the platform on a single-node cluster. Evaluation covers 26 requirements, of which two are verified functionally, five are instrumented with their targets left unmeasured, and 19 are verified structurally. The retraining chain fires when injected drift crosses the configured threshold, yielding a Population Stability Index of 18.107 against a threshold of 0.2. A cost comparison across three provisioning options places the open source design below fully managed services over three years under an assumed labour allocation. Load characterisation under real traffic remains open, on one node and across a cluster.
\end{abstract}

\begin{IEEEkeywords}
DataOps, drift detection, feature store, GitOps, Kubernetes, MLOps, open source
\end{IEEEkeywords}
